\chapter{\MakeUppercase{Technical Specification}}

\section{Requirements}

\subsection{Functional requirements}

\begin{enumerate}
    \item I wanted the system to take any medical question in English and return a direct text answer.
    \item I forced the pipeline to pull relevant facts from the knowledge base before it even starts writing.
    \item I made sure every response points back to the specific documents that contributed to the answer.
    \item I added a confidence score calculation that attaches to every response.
    \item I built a hallucination gate to verify that my generated outputs actually match the retrieved context.
    \item I included safety guardrails to detect and redirect any queries about emergencies, drug dosing, or self-harm.
    \item I set up a \ac{rest} \ac{api} with endpoints for standard queries (\texttt{/ask}), streaming, and health checks.
    \item I implemented query caching so I can serve repeat questions without re-running the whole process.
    \item I designed the interface to show the answer, trust levels, source lists, and speed metrics all at once.
\end{enumerate}

\subsection{Non-functional requirements}

\begin{enumerate}
    \item I ensured the entire pipeline runs locally on a \ac{cpu} with no internet connection required.
    \item I prioritised privacy by making sure the system doesn't log or store any user queries.
    \item I kept container startup time under five minutes on my test hardware.
    \item I targeted a warm-path \ac{api} latency of less than 200 seconds for TinyLlama on a \ac{cpu}.
    \item I wrote a suite of 250 automated tests that I use to verify unit, integration, and safety logic.
    \item I used Docker to ensure the runtime behavior is identical across different host machines.
    \item I secured the \ac{api} with keys and separate scopes for basic read and administrative access.
\end{enumerate}

\subsection{Safety and guardrail requirements}

\begin{enumerate}
    \item Questions matching emergency patterns --- chest pain, difficulty breathing, suicidal ideation, stroke symptoms --- shall receive a redirect to emergency services, not a generated answer.
    \item Questions requesting specific prescription dosages shall receive a redirect to consult a physician.
    \item All responses shall include a disclaimer that the system is not a substitute for professional medical advice.
    \item The system shall use an aggressive stop-word list to prevent the generation model from producing exam-format MCQ patterns (e.g., ``A) \ldots B) \ldots C) \ldots'') when the underlying dataset contains such formats.
\end{enumerate}

\section{Feasibility study}

\subsection{Technical feasibility}

I calculated that the 1.1B-parameter TinyLlama model, when loaded in bfloat16, occupies about 2.2 GB of \ac{ram}. My test machine had 16 GB available, which left me with plenty of headroom once I factored in the 2--4 GB needed for retrieval (the \ac{bm25} index and ChromaDB vector store). Although I noticed a warm-path latency of 106 seconds, which is too slow for real-time consumer apps, I found it perfectly acceptable for my research prototype or for use in an asynchronous queue. I also tested the BioMistral-7B model in 4-bit \ac{gguf} quantisation; it requires roughly 4 GB and resulted in a higher mean latency of 143 seconds.

\subsection{Data feasibility}

I gathered three main datasets from public sources under open research licences. I started with ChatDoctor-HealthCareMagic, which provided me with 226,000 real patient-doctor conversation pairs. Next, I integrated PubMedQA, where I found 1,000 expert-annotated pairs and another 61,000 unlabelled ones. I finished the collection with 187,000 postgraduate entrance exam questions from MedMCQA. After I filtered, deduplicated, and applied sentence-boundary chunking, my combined knowledge base reached a total of 505,584 documents.

\subsection{Operational feasibility}

I observed that starting the container takes about 2--3 minutes on my test bed, mainly because loading TinyLlama in bfloat16 on a \ac{cpu} is a slow process. To mitigate this, I included a \texttt{SKIP\_STARTUP\_WARMUP=true} flag that delays the model load until the first query arrives, letting me iterate faster at the cost of a slower first response.

\section{System specification}

\subsection{Hardware specification}

\begin{table}[H]
\centering
\caption{Hardware used for local development and evaluation}
\renewcommand{\arraystretch}{1.1}
\begin{tabular}{|l|l|}
\hline
\textbf{Component} & \textbf{Specification} \\
\hline
Processor & x86-64 \ac{cpu} (no \ac{gpu}) \\
RAM & 16 GB \\
Storage & \ac{ssd}, 500 GB \\
Operating system & Ubuntu 22.04 / Ubuntu 24.04 \\
Training (Colab) & \ac{tpu} v5e (knowledge base build), T4 \ac{gpu} (BioMistral eval) \\
\hline
\end{tabular}
\end{table}

\subsection{Software specification}

\begin{table}[H]
\centering
\caption{Core software components and versions}
\renewcommand{\arraystretch}{1.1}
\begin{tabular}{|l|l|l|}
\hline
\textbf{Component} & \textbf{Version} & \textbf{Purpose} \\
\hline
Python & 3.11 & Runtime \\
FastAPI & 0.110+ & \ac{rest} \ac{api} framework \\
LangChain & 0.1+ & Pipeline orchestration \\
LangGraph & 0.1+ & Stateful pipeline graphs \\
ChromaDB & 0.4+ & Vector store \\
sentence-transformers & 2.7+ & Dense embeddings \\
rank-bm25 & 0.2.2 & Sparse retrieval \\
llama-cpp-python & 0.2+ & \ac{gguf} inference (BioMistral) \\
transformers & 4.40+ & TinyLlama inference \\
peft & 0.10+ & QLoRA adapter loading \\
React (assistant-ui) & 18+ & Frontend \\
Docker & 24+ & Containerisation \\
pytest & 8+ & Testing \\
\hline
\end{tabular}
\end{table}

\subsection{Knowledge base specification}

\begin{table}[H]
\centering
\caption{Knowledge base v2 composition}
\renewcommand{\arraystretch}{1.1}
\resizebox{\textwidth}{!}{
\begin{tabular}{|l|l|r|}
\hline
\textbf{Dataset} & \textbf{Type} & \textbf{Docs after chunking} \\
\hline
ChatDoctor-HealthCareMagic & Patient-doctor QA pairs & 226,000 \\
PubMedQA & Biomedical research abstracts & 62,000 \\
MedMCQA & Medical MCQ pairs & 217,000 \\
\hline
Total & -- & 505,584 \\
\hline
\end{tabular}}
\end{table}

For my dense retrieval module, I relied on the \texttt{all-MiniLM-L6-v2} model from SentenceTransformers to generate 384-dimensional vectors. I chose to store these in ChromaDB using cosine similarity to ensure matching accuracy. For the sparse retrieval component, I built a \ac{bm25} index from the same document set; I configured the application to load this index from a pickle file at \path{data/knowledge_base_v2/bm25_index.pkl} during the startup phase.
